% ===========================================================================
% Section 4: Dataset and Preprocessing (~1.5 pages)
% ===========================================================================

% --- 4.1 Dataset ---
\subsection{Dataset}
\label{sec:dataset:data}

% TODO: Describe the Polaris dataset
% - 1.37M Darshan logs from ALCF Polaris (Apr 2024 -- Feb 2026)
% - Anonymized: uid/exe hashed, mount points preserved
% - Storage: Eagle/Grand Lustre (160 OSTs, 40 MDTs, 100 PiB, 650 GiB/s)
% - Cite: \cite{snyder2025darshan}

% Table: Dataset summary
\begin{table}[t]
\centering
\caption{Dataset Summary}
\label{tab:dataset}
\begin{tabular}{lr}
\toprule
\textbf{Property} & \textbf{Value} \\
\midrule
Source system & ALCF Polaris \\
Collection period & Apr 2024 -- Feb 2026 \\
Raw Darshan logs & 1,397,216 \\
After cleaning & 131,151 \\
Retention rate & 9.4\% \\
Total features & 186 \\
\midrule
\multicolumn{2}{l}{\textit{Feature breakdown:}} \\
~~Raw counters & 136 \\
~~Indicators & 8 \\
~~Derived ratios & 40 \\
~~Metadata & 2 \\
\midrule
Train / Val / Test & 91,807 / 19,672 / 19,672 \\
Split method & Temporal (70/15/15) \\
\bottomrule
\end{tabular}
\end{table}

% --- 4.2 Feature Extraction ---
\subsection{Feature Extraction}
\label{sec:dataset:features}

% TODO: Describe feature extraction
% - 186 raw features from Darshan POSIX/MPI-IO/STDIO modules
%   - 136 raw counters (aggregated across files per job)
%   - 8 indicator features (has_mpiio, has_stdio, etc.)
%   - 39 derived ratios and metrics
%   - 3 metadata (nprocs, duration, total_bytes)
% - Aggregation rules: SUM, MAX, MIN_NONZERO, LAST_VALUE, TOP-4 MERGE
% - 157 ML-ready features after exclusion of timestamps, constants, hash values

% --- 4.3 Preprocessing Pipeline ---
\subsection{Preprocessing Pipeline}
\label{sec:dataset:preprocessing}

% TODO: Describe preprocessing stages
% Stage 1: Extraction (1,397,293 files)
% Stage 2: Cleaning (-> 131,151 rows after filtering)
%   - min_duration >= 1s, min_total_bytes >= 1024, min_io_ops >= 1
% Stage 3: Feature engineering (39 derived features)
% Stage 4: EDA and feature exclusion (186 -> 157 features)
% Stage 5: Group-specific normalization (14 feature groups)

% Figure: Data characterization (3 panels)
% \begin{figure}[t]
%   \centering
%   \includegraphics[width=\columnwidth]{fig_data_characterization.pdf}
%   \caption{Dataset characterization: (a) I/O volume distribution,
%   (b) feature sparsity, (c) normalization effect.}
%   \label{fig:data_char}
% \end{figure}

% Table: Normalization methods
\begin{table}[t]
\centering
\caption{Group-Specific Normalization Strategy}
\label{tab:normalization}
\begin{tabular}{lrl}
\toprule
\textbf{Feature Group} & \textbf{\#} & \textbf{Method} \\
\midrule
    Volume & 12 & $\log(1{+}x)$ + RobustScaler \\
    Count & 39 & $\log(1{+}x)$ + RobustScaler \\
    Histogram & 20 & $\log(1{+}x)$ \\
    Top4 & 24 & $\log(1{+}x)$ \\
    Timing & 25 & $\log(1{+}x)$ + RobustScaler \\
    Ratio Unbounded & 14 & $\log(1{+}x)$ \\
    Derived Absolute & 3 & $\log(1{+}x)$ \\
    Metadata & 2 & $\log(1{+}x)$ \\
    Conditional Size & 2 & $\log(1{+}x)$ \\
    Ratio & 23 & none \\
    Indicator & 8 & none \\
    Timestamp & 8 & none \\
    Categorical & 4 & none \\
    Rank Id & 2 & none \\
\bottomrule
\end{tabular}
\end{table}

% --- 4.4 Temporal Data Split ---
\subsection{Temporal Data Split}
\label{sec:dataset:split}

% TODO: Describe temporal split
% - 70/15/15 split by _start_time (chronological order)
% - Train: 91,807 (earliest), Val: 19,672, Test: 19,672 (latest)
% - Justified by: TabReD (ICLR'25), Kim et al. (J.BigData'23)
% - Prevents data leakage from temporal autocorrelation
% - Polaris had 5+ system upgrades in collection period

% Figure: Temporal split
% \begin{figure}[t]
%   \centering
%   \includegraphics[width=\columnwidth]{fig_temporal_split.pdf}
%   \caption{Temporal train/val/test split showing job volume over time.
%   Vertical lines indicate split boundaries.}
%   \label{fig:temporal_split}
% \end{figure}
